\chapter{External experimental validation}
\label{ch:external}

\begin{saybox}[title={What to say at the start of this section}]
This is the chapter that answers ``does any of this work on real cheese?''
The honest answer is: on the evidence in this repository, not yet, and the
evaluation that exists is small and flagged as out-of-support. Present it
plainly --- a panel will respect a stated negative result far more than an
elided one.
\end{saybox}

\section{What external evidence exists}

Three external files are present in the repository:

\begin{center}\small
\begin{tabular}{lrp{6.0cm}}
\toprule
File & rows & Role \\
\midrule
\texttt{CHEESE\_100\_EXTERNAL\_LITERATURE\_TEST\_CASES.xlsx} & 100 &
Candidate literature cases with reported shelf-life outcomes \\
\texttt{CHEESE\_319\_REAL\_EXTERNAL\_STRESS\_TEST\_CASES.xlsx} & 319 &
Larger stress-test case set \\
\texttt{reports/v7\_external\_test\_predictions.csv} & \textbf{21} &
V7 predictions for the subset of the 100 that the specialists could
actually predict \\
\bottomrule
\end{tabular}
\end{center}

Of the $100$ candidate literature cases, only $21$ were predictable by the
V7 specialists; the remainder were rejected by the routing and support
checks before a prediction was produced. The $21$ come from $6$ distinct
source citations.

\section{The external result}

\begin{table}[htbp]\centering\small
\caption{V7 external literature evaluation, all 21 predictable cases.
Computed directly from \texttt{reports/v7\_external\_test\_predictions.csv}.}
\label{tab:external}
\begin{tabular}{lrrr}
\toprule
Metric & Value \\
\midrule
Cases predicted & 21 (of 100 candidates) \\
Distinct source citations & 6 \\
Best model selected & LightGBM for all 21 \\
External $\rsq$ & \textbf{0.172} \\
MAE & \textbf{39.4 days} \\
Mean signed error & $\mathbf{-30.0}$ days (systematic under-prediction) \\
Cases flagged \texttt{unsupported\_categorical} & \textbf{21 of 21 ($100\%$)} \\
\bottomrule
\end{tabular}
\end{table}

\begin{findingbox}[title={The external result, stated without softening}]
On the only real-world evaluation available in this repository, the V7
system achieves $\rsq = 0.172$ with a mean absolute error of $39.4$ days
and a systematic $30$-day under-prediction.

Against a leave-rule-out pooled out-of-fold $\rsq$ between $0.67$ and
$0.96$ on synthetic data, this is the quantitative statement of the
sim-to-real gap. It is the single most important number in the report for
any question about deployment.
\end{findingbox}

\subsection{Individual cases}

\begin{center}\small
\begin{tabular}{lrrr}
\toprule
Product & true (d) & predicted (d) & error (d) \\
\midrule
Mozzarella & 13 & 10.9 & $-2.1$ \\
Mozzarella & 63 & 12.5 & $-50.5$ \\
Mozzarella & 45 & 12.4 & $-32.6$ \\
Mozzarella & 16 & 11.8 & $-4.2$ \\
High-moisture Mozzarella & 8 & 11.2 & $+3.2$ \\
High-moisture Mozzarella & 4 & 10.9 & $+6.9$ \\
Fresh acid-coagulated cheese & 9 & 13.7 & $+4.7$ \\
Processed Gouda spread & 28 & 43.1 & $+15.1$ \\
Processed Gouda spread & 4 & 43.1 & $+39.1$ \\
Processed Gouda spread & 49 & 53.0 & $+4.0$ \\
Montasio & 60 & 78.4 & $+18.4$ \\
Montasio & 120 & 78.4 & $-41.6$ \\
Provolone (aged 6 months) & 100 & 59.5 & $-40.5$ \\
Provolone (aged 6 months) & 175 & 64.6 & $-110.4$ \\
Provolone (aged 6 months) & \textbf{280} & \textbf{59.8} & $\mathbf{-220.2}$ \\
Provolone (aged 6 months) & 190 & 60.1 & $-129.9$ \\
\bottomrule
\end{tabular}
\end{center}

\noindent (Sixteen of the twenty-one cases shown; the remaining five are
further Mozzarella and Montasio replicates with the same pattern.)

Two failure modes are visible and they are different from each other:

\paragraph{Insensitivity within a product.} Four Mozzarella cases with true
values of $13$, $63$, $45$ and $16$ days all receive predictions between
$10.9$ and $12.5$ days. The model is not merely inaccurate; it is nearly
constant across cases that differ by a factor of five in reality. The same
occurs for Montasio ($60$--$120$ days true, all predicted $78.4$) and
Provolone ($100$--$280$ days true, all predicted $59.5$--$64.6$).

\paragraph{Systematic under-prediction on aged hard cheese.} The five
Provolone cases are under-predicted by $40$ to $220$ days. This is
\cref{eq:rf-bound} again, now against real data: the synthetic hard-cheese
training distribution does not extend to the shelf lives that aged
Provolone actually achieves.

\begin{findingbox}[title={A historical note worth raising at the meeting}]
Project documentation for the pre-specialist v1 model (2026-08-13) records
the \emph{opposite} bias on the same external cases: Provolone true $45$ d
predicted $126$ d ($+180\%$), Fresh Mozzarella true $4$ d predicted $5.9$ d
($+48\%$) --- systematic \emph{over}-prediction, with an external MAE of
$50$--$100$ days.

The V7 system now under-predicts by $30$ days on average. The sign of the
bias has flipped across dataset revisions while the magnitude has roughly
halved. That is consistent with successive recalibrations moving the
synthetic target distribution without ever anchoring it to real data ---
and it means the external gap has been a known, tracked property of this
project since its first model version, not a recent discovery.
\end{findingbox}

\section{The independence question}

The audit request asked whether the external cases are genuinely
independent of the synthetic generation rules. The honest answer is that
\textbf{this cannot be determined from this repository}.

\begin{limitbox}[title={Independence is unverifiable here}]
The external cases carry \texttt{citation}, \texttt{source\_doi} and
\texttt{publication\_year}. The generation rules carry only composite
identifiers such as \texttt{NATAMYCIN\_REVIEW} or \texttt{RICOTTA\_MAP},
which reference source \emph{components} without bibliographic metadata.
Without the generator's source-to-rule mapping (absent, per
\Cref{sec:limit-evidence}), there is no way to test whether a publication
behind an external test case also informed a generation rule.

Component names such as \texttt{HMMC\_NAT} (high-moisture Mozzarella,
natamycin) and the presence of high-moisture Mozzarella among the external
cases make \emph{some} overlap plausible. Plausible is not established.
This should be raised with whoever maintains the generation pipeline, and
it is the single highest-value piece of missing evidence in the whole
audit.
\end{limitbox}

\subsection{The support flag}

Every one of the $21$ predictions carries
\texttt{support\_level = unsupported\_categorical}: the serving layer
itself recorded that each case contained a categorical value outside the
training support. The system flagged its own predictions as
out-of-distribution and produced them anyway for evaluation purposes.

This is, in its way, a positive finding about the engineering: the support
check works and it fired correctly on all $21$ cases. It also means the
external $\rsq = 0.172$ should be read as ``performance when
extrapolating beyond categorical support,'' which is a harder test than a
well-supported case would face --- but since no externally-validated case
in the repository \emph{is} well-supported, no easier number exists.

\section{What would be needed for a real-world claim}

\begin{enumerate}[leftmargin=1.5em, itemsep=2pt]
\item A set of real cases that fall \emph{inside} the categorical support
  of the trained specialists, so that the evaluation measures prediction
  rather than extrapolation.
\item Verified independence between those cases' source publications and
  the generation rules --- which requires the generator's source mapping.
\item Enough cases per specialist to give the estimate a usable interval.
  Twenty-one cases across three categories and six citations does not.
\item Reporting of the error in days alongside $\rsq$, since with $n=21$
  and a wide true range the $\rsq$ is dominated by a handful of Provolone
  cases.
\end{enumerate}

Until those exist, the defensible position is the one this report takes
throughout: the system has demonstrated cross-rule transfer on synthetic
data, and has not demonstrated real-world accuracy.
